
\documentclass{article}
\usepackage{amsmath}


\usepackage{graphicx,comment,framed}
\usepackage{roundbox}
\usepackage{fancybox}
\usepackage{tikz}
\usepackage{color}
\usepackage[hidelinks]{hyperref}
\usepackage{framed}
\usepackage{amsthm,amssymb,amsmath}
%\usepackage[colorlinks,linkcolor=blue,citecolor=blue]{hyperref}
\definecolor{shadecolor}{cmyk}{0,0,0,0}
\usepackage{listings}
\usepackage{xepersian}
\usepackage[noend]{algpseudocode}


%--------------------- page settings ----------------------

\settextfont[Scale=1.1]{XB Niloofar}
\setdigitfont[Scale=1.1]{XB Niloofar}
\defpersianfont\sayeh[Scale=1.1]{XB Niloofar}
\addtolength{\textheight}{3.2cm}
\addtolength{\topmargin}{-22mm}
\addtolength{\textwidth}{3cm}
\addtolength{\oddsidemargin}{-1.5cm}


%------------------------ Environments ------------------------------------

\newtheorem{قضیه}{قضیه}
\newtheorem{لم}{لم}
\newtheorem{مشاهده}{مشاهده}
\newtheorem{تعریف}{تعریف}


%-------------------------- Notations ------------------------------------

\newcommand{\IR}{\ensuremath{\mathbb{R}}} 
\newcommand{\IZ}{\ensuremath{\mathbb{Z}}} 
\newcommand{\IN}{\ensuremath{\mathbb{N}}} 
\newcommand{\IS}{\ensuremath{\mathbb{S}}} 
\newcommand{\IC}{\ensuremath{\mathbb{C}}} 
\newcommand{\IB}{\ensuremath{\mathbb{B}}} 

\newcommand{\bR}{\mathbb{R}}
\newcommand{\cB}{\mathcal{B}}
\newcommand{\cO}{\mathcal{O}}
\newcommand{\cG}{\mathcal{G}}
\newcommand{\rM}{\mathrm{M}}
\newcommand{\rC}{\mathrm{C}}
\newcommand{\rV}{\mathrm{V}}

\newcommand{\lee}{\leqslant}
\newcommand{\gee}{\geqslant}
\newcommand{\ceil}[1]{{\left\lceil{#1}\right\rceil}}
\newcommand{\floor}[1]{{\left\lfloor{#1}\right\rfloor}}
\newcommand{\prob}[1]{{\mbox{\tt Pr}[#1]}}
\newcommand{\set}[1]{{\{ #1 \}}}
\newcommand{\seq}[1]{{\left< #1 \right>}}
\newcommand{\provided}{\,|\,}
\newcommand{\poly}{\mbox{\rm poly}}
\newcommand{\polylog}{\mbox{\rm \scriptsize polylog}\,}
\newcommand{\comb}[2] {\left(\!\!\begin{array}{c}{#1}\\{#2}\end{array}\!\!\right)}




\newcounter{probcnt}
\newcommand{\مسئله}[1]{\stepcounter{probcnt}\bigskip\bigskip{
 	\large \bf مسئله‌ی \arabic{probcnt}$\mbox{\bf{.}}$ \ #1} \bigskip}

\newcommand{\fqed}[1]{\leavevmode\unskip\nobreak\quad\hspace*{\fill}{\ensuremath{#1}}}

\newenvironment{اثبات}
	{\begin{trivlist}\item[\hskip\labelsep{\em اثبات.}]}
	{\fqed{\square}\end{trivlist}}

\newenvironment{حل}
	{\begin{trivlist}\item[\hskip\labelsep{\bf حل.}]}
	{\fqed{\blacktriangleright}\end{trivlist}}

\ifdefined\hidesols
	\newsavebox{\trashcan} % uncomment the following line to hide solutions
	\renewenvironment{حل}{\begin{latin}\begin{lrbox}{\trashcan}}{\end{lrbox}\end{latin}}
\fi


%------------------------- Header -----------------------------

%------------------------- Header -----------------------------

\newcommand{\سربرگ}[3]{
	\parindent=0em
	\begin{shaded}
		
		\rightline{ 
			\makebox[6em][c]{
				\includegraphics[height=1.6cm]{aut.png}
		}}
		\vspace{-.2em}
		{\scriptsize\bf دانشکده‌ی مهندسی کامپیوتر}
		\hfill {\small
			مدرس: دکتر امیرمزلقانی  \ 
		}\\[-5em]
		\leftline{\hfill\Large\bf 
			کاربرد های جبر خطی
		}\\[.7em]
		\leftline{\hfill\bf 
			نیم‌سال دوم ۹۶-۹۷
		}\\[1.7em]
		\hrule height .12em
		
		\normalsize
		\vspace{1mm} #1
		\hfill \small  #3
		\vspace{1mm} 
		\hrule height .1em
		
		\vspace{-0.5em} 
		\hfill {\sayeh\large #2} \hfill
	\end{shaded}
	\begin{large}

توجه:
\end{large}
\\
\\
\begin{itemize}
\item
این تمرین از مباحث مربوط به فصل5 ،6و7 طراحی شده است که شامل 15سوال اجباری و 4 سوال امتیازی است که نمره سوال های امتیازی فقط به نمرات تمرین شما کمک می کند.
\item 
در برخی سوالات چند قسمتی از شما خواسته شده است به تعدادی از قسمت ها به اختیار پاسخ دهید ،پاسخگویی به بیش از مقدار تعیین شده نمره اضافی ندارد ،از بقیه قسمت ها می توانید برای تمرین بیشتر استفاده کنید.
\item 
در پایان نیز تعدادی سوال برای تمرین بیشتر در نظر گرفته شده است ،که به حل آن ها نمره ای تعلق نمی گیرد صرفا برای تمرین توسط خود شما در نظر گرفته شده است. 
\item 
اگه سوالی داشتین از طریق
\begin{latin}
\href{mailto:aut.la2018@gmail.com}{\[aut.la2018@gmail.com\]}
\end{latin}
 حتما بپرسید.

\item پاسخ های تمرین را در قالب یک فایل به صورت الگوی زیر آپلود کنید.
\begin{latin}
9531000\_steve\_McManaman\_HW4.pdf
\end{latin}
\item \color{red}
مهلت تحویل جمعه  1تیر 1397 ساعت 23:54:59
\end{itemize}

}




%\settextfont{Yas}
\begin{document}
\سربرگ{تمرین آشنایی با متلب}{}{}
\clearpage
\مسئله{متغیر های اسکالر}

متغیر های زیر را ایجاد کنید:
\lr{
\begin{enumerate}
	\item $a = 10$
	\item $b = 1023 * 2.5$
	\item $c = 2+3j$ \text{عدد موهومی در اعداد مختلط است}
	\item $d = e^{j2pi}$  (use \textbf{exp}, \textbf{pi})
\end{enumerate}
}
j) عدد موهومی در اعداد مختلط است)

\مسئله{متغیر های برداری}
\begin{latin}
	\begin{enumerate}
		\item aٰVec = [3.14 15 9 26]
		
		\item bVec = $\begin{bmatrix}
			2.17 \\
			8 \\
			28 \\
			182
		\end{bmatrix}$
		
		\item cVec = [5 4.8 4.6  \dots  -4.8 -5] 
		
		(all the numbers from 5 to -5 in increments of -0.2)
		
		\item dVec = [$10^0$ $10^{0.01}$ $10^{0.02}$ \dots $10^{0.99}$ $10$ ] 
		
		(logarithmically spaced numbers between 1and 10, use \textbf{logspace}, make sure you get the length right!)
		
		\item eVec = Hello 
		
		( eVec is a string, which is a vector of characters)
	\end{enumerate}
\end{latin}

\مسئله{متغییر های ماتریسی}

\begin{latin}
	\begin{enumerate}
		\item
		$
		 aٰMat = \begin{bmatrix}
			2  & \dots &2 \\
			\vdots  & \ddots & \vdots \\
			2  & \dots &2 \\
		\end{bmatrix}
		$ 
		
		(a 9x9 matrix full of 2’s (use \textbf{ones} or \textbf{zeros}))
		\item
		$
		bMat = \begin{bmatrix}
			1  & \dots &\dots &\dots &0 \\
			\vdots  & \ddots &0&\ddots&\vdots \\
			\vdots  & 0&5&0&\vdots\\
			\vdots  & \ddots&0& \ddots&\vdots\\
			0  & \dots & \dots& \dots &1 \\
		\end{bmatrix}
		$ 
		
		(a 9x9 matrix of all zeros, but with the values [1 2 3 4 5 4 3 2 1] on the main diagonal (use \textbf{zeros},\textbf{diag}).)
		
		\item
		$
		cMat = \begin{bmatrix}
		1  & 11 &\dots &\dots &91 \\
		2  & 12 &\ddots&\ddots&92 \\
		\vdots  & \ddots&\ddots&\ddots&\vdots\\
		9  & 19&\ddots& \ddots&99\\
		10  & 20 & \dots& \dots &100 \\
		\end{bmatrix}
		$ 
		
		(a 10x10 matrix where the vector 1:100 runs down the columns (use \textbf{reshape}).)

	\end{enumerate}
\end{latin}
\مسئله{معادلات اسکالر}
\begin{latin}
	\begin{enumerate}
		\item x = {\LARGE${\frac{1}{1 + e^{-(a-15)/6}}}$}
		\item y = $(\sqrt{a} + \sqrt[21]{b}) ^{pi}$
	\end{enumerate}
\end{latin}

\مسئله{معادلات برداری}

در این قسمت از معادلات بخش 2 و عملگرهای عنصر نهاد
(\lr{element wise})
 ( مثل $.*$ و $./$ و...) استفاده کنید. (برای مثال، عناصر xVec حاصل محاسبه فرمول فوق بر روی تک تک عناصر cVec هستند)
\begin{latin}
	\begin{enumerate}
		\item xVec = {\LARGE${\frac{1}{\sqrt{2pi * 2.5}}}e^{{-cVec^2}/2}$}
		\item yVec = $\sqrt{(aVec)^T + bVec^2}$
	\end{enumerate}
\end{latin}
\مسئله{معادلات ماتریسی}

در این قسمت از عملگرهای ماتریسی استفاده کنید.
\begin{latin}
	\begin{enumerate}
		\item xMat = $(aVec * bVec) * aVec^T$
	\end{enumerate}
\end{latin}
\مسئله{ماتریس ستونی ای تشکیل دهید که میانگین هر سطر cMat را در خود ذخیره کند. همچنین میانگین کلی ماتریس را نیز پیدا کرده و در متغیری ذخیره و نمایش دهید. (‌از تابع \textbf{mean} استفاده کنید)}

\end{document}